应用 Lemma~\ref{lem:flat-boundary-local-regularity} 之前, 需要证明变换后的微分算子一致椭圆.

\begin{Lemma}
\label{lem:boundary-transformed-uniform-ellipticity}
由
\begin{equation}\label{eq:boundary-transformed-second-order-operator}
\widetilde L(\bm{y},D)u:=\sum_{k,l=1}^n(\overline a_{kl}(\bm{y})v_{y_l})_{y_k}+\sum_{k=1}^n\overline b_k(\bm{y})v_{y_k}+\overline c(\bm{y})v
\end{equation}
定义的算子 $\widetilde L$ 在 $\Omega^{\prime\prime}$ 中一致椭圆.
\end{Lemma}

\begingroup
\renewcommand{\endProof}{\par\nobreak\hfill$\blacksquare$\par\vspace{0.1cm}}
\begin{Proof}
必须证明存在常数 $\widetilde\theta>0$, 使得
\begin{equation}\label{eq:boundary-transformed-principal-ellipticity-bound}
-\sum_{k,l=1}^n\overline a_{kl}(\bm{y})\xi_k\xi_l\geq\widetilde\theta|\bm{\xi}|^2
\end{equation}
对每个 $\bm{\xi}\in\mathbb{R}^n$ 和每个 $\bm{y}\in\Omega^{\prime\prime}$ 都成立.

对任意 $\bm{\xi}\in\mathbb{R}^n$, 令 $\bm{\eta}:=\bm{A}\bm{\xi}$, 其中\[\begin{aligned}
\bm{A}(\bm{y})&:=\nabla\Psi(\Phi(\bm{y}))^T\\
&=\begin{pmatrix}
1&0&\cdots&0&-\psi_{x_1}\\
0&1&\cdots&0&-\psi_{x_2}\\
\vdots&\vdots&\ddots&\vdots&\vdots\\
0&0&\cdots&1&-\psi_{x_{n-1}}\\
0&0&\cdots&0&1
\end{pmatrix}.
\end{aligned}\]注意 $\bm{A}(\bm{y})$ 可逆. 令
\begin{equation}\label{eq:boundary-inverse-jacobian-uniform-bound}
\widetilde C:=\sup_{\bm{y}\in\Omega^{\prime\prime}}|\bm{A}^{-1}(\bm{y})|.
\end{equation}
则
\begin{equation}\label{eq:boundary-jacobian-vector-norm-comparison}
|\bm{\xi}|\leq\widetilde C|\bm{\eta}|.
\end{equation}
现在利用 \eqref{eq:boundary-jacobian-vector-norm-comparison} 和 $L$ 的一致椭圆性得到\[\begin{aligned}
-\sum_{k,l=1}^n\overline a_{kl}\xi_k\xi_l
&=-\sum_{i,j=1}^na_{ij}(\Phi(\bm{y}))\sum_{k,l=1}^n\left(\frac{\partial\Psi_k}{\partial x_i}(\Phi(\bm{y}))\xi_k\right)\left(\frac{\partial\Psi_l}{\partial x_j}(\Phi(\bm{y}))\xi_l\right)\\
&=-\sum_{i,j=1}^na_{ij}(\phi(\bm{y}))\eta_i\eta_j\\
&\geq\theta|\bm{\eta}|^2\\
&\geq\frac{\theta}{\widetilde C^2}|\bm{\xi}|^2.
\end{aligned}\]因此, $\widetilde L$ 以常数 $\widetilde\theta:=\theta/\widetilde C^2$ 一致椭圆.
\end{Proof}
\endgroup
